\PassOptionsToPackage{dvipsnames,svgnames,table}{xcolor}
\documentclass[10pt]{article}
\usepackage[a4paper,margin=16mm]{geometry}
\usepackage[no-math]{fontspec}\setmainfont{Latin Modern Roman}
\usepackage{amsmath,amssymb,amsthm,mathtools,graphicx,xcolor,hyperref,xurl}
\usepackage[shortlabels]{enumitem}
\setlength{\emergencystretch}{4em}
\newtheorem{theorem}{Theorem}[section]
\newtheorem{thm}[theorem]{Theorem}
\newtheorem{lemma}[theorem]{Lemma}
\newtheorem{lemm}[theorem]{Lemma}
\newtheorem{proposition}[theorem]{Proposition}
\newtheorem{prop}[theorem]{Proposition}
\newtheorem{corollary}[theorem]{Corollary}
\newtheorem{coro}[theorem]{Corollary}
\newtheorem{conjecture}[theorem]{Conjecture}
\theoremstyle{definition}
\newtheorem{definition}[theorem]{Definition}
\newtheorem{defi}[theorem]{Definition}
\newtheorem{defn}[theorem]{Definition}
\newtheorem{remark}[theorem]{Remark}
\newtheorem{rema}[theorem]{Remark}
\newtheorem{example}[theorem]{Example}
\newtheorem{exam}[theorem]{Example}
\makeatletter\let\c@cdrthm\c@theorem\let\thecdrthm\thetheorem\let\p@cdrthm\p@theorem\let\theHcdrthm\theHtheorem\makeatother
\providecommand{\eg}{e.g.\ }
\providecommand{\ie}{i.e.\ }
\providecommand{\etc}{etc.\ }
\providecommand{\cf}{cf.\ }
\providecommand{\longthanks}[1]{\par\smallskip\noindent #1\par\smallskip}
\DeclareMathOperator{\card}{card}
\DeclareMathOperator{\Supp}{Supp}

\newcommand{\C}{{\mathbb C}}
\newcommand{\N}{{\mathbb N}}
\newcommand{\R}{{\mathbb R}}
\newcommand{\Q}{{\mathbb Q}}
\newcommand{\Z}{{\mathbb Z}}
\newcommand{\F}{{\mathbb F}}


%%%%%%%%%%%%%%%%%%%%%%%%%%%%%%%%%%%%%%%%%%%%%%%%%%%%%%%%%%%%%%%%%




%%%%%%%%% a placer avant \begin{document}


%%%%%%%%%%%%%%%%%%%%%%%%%%%%%%%%%%%%%
\newcommand*{\mk}{\mkern -1mu}
\newcommand*{\Mk}{\mkern -2mu}
\newcommand*{\mK}{\mkern 1mu}
\newcommand*{\MK}{\mkern 2mu}

\hypersetup{urlcolor=purple, linkcolor=blue, citecolor=red}


\newcommand*{\romanenumi}{\renewcommand*{\theenumi}{\roman{enumi}}}
\newcommand*{\Romanenumi}{\renewcommand*{\theenumi}{\Roman{enumi}}}
\newcommand*{\alphenumi}{\renewcommand*{\theenumi}{\alph{enumi}}}
\newcommand*{\Alphenumi}{\renewcommand*{\theenumi}{\Alph{enumi}}}
%%%%%%%%%%%%%%%%%%%%%%%%%%%%%%%%%%



\begin{document}
\section*{1313. Degree and quotient-double-coset criterion for bi-anchored thin chains in standard complex table algebras}
Blau, Harvey I.. \emph{On residually thin and nilpotent table algebras, fusion rings, and association schemes}. Algebraic Combinatorics 5 (2022), publisher version, DOI 10.5802/alco.194. \url{https://alco.centre-mersenne.org/articles/10.5802/alco.194/}. Exact archived official native/PDF version. Native and publisher PDF individually CC BY4.0, \url{https://creativecommons.org/licenses/by/4.0/}.
\paragraph{Selected problem (editorial).} Prove exactly Theorem1.7: for a finite-dimensional complex standard table algebra and a strictly nested bi-anchored chain, thinness is equivalent to the printed degree/product-of-successive-double-coset-cardinalities formula for every1<=i<=n and b in C\_(i-1) minus C\_i. One bidirectional criterion is selected. The complete new chain-factorization and both degree-constancy quotient inductions are included. Later divisibility, p-standard, nilpotence, and counterexample assertions are preserved as uncounted author context.
\paragraph{Ordinary prerequisites (editorial).} Frobenius--Perron degree-map foundations, standard table-algebra quotient/homomorphism theory and the closed-subset correspondence are explicitly cited original prerequisites. Lemma3.3 and Theorem3.4 and both thinness implications remain new core, never granted.
\paragraph{Difficulty (editorial).} H1. The target recovers every thin quotient in a nested standard-table-algebra chain from a degree formula, through a new double-coset degree-constancy criterion and recursive quotient-chain reconstruction. Even after the named quotient and homomorphism prerequisites, the bidirectional induction and exact multiplicative transfer remain substantive proof obligations beyond routine qualification-exam arguments. No later integrality or nilpotence result contributes to this rating.
\paragraph{Literal source and selected scope (editorial).} The exact author proof retains its literal i=0 phrases at952/956/958, while the selected original statement and definition of m at294--307 use1<=i<=n. No bound or C\_(i-1) edit is supplied. The vacuous b in B minus C\_0 language in the supporting criterion is preserved. Root source decision remains separate.
\paragraph{Formatting (editorial).} Standalone article scaffold replaces publisher administrative class/front matter only. Original active macro definitions, mathematics, indices, proofs and percent guards within every selected block are unchanged. No publisher class, private preamble or external software acquisition.
\clearpage
\section{Introduction}\label{sect1}
We explore two related aspects of
some important algebraic and combinatorial structures: namely, the
properties called \emph{residual thinness} and \emph{nilpotence}. The
former concept has been studied in the framework of association schemes
by Zieschang~\cite{Z1, Z2} and Hanaki and Shimabukuro~\cite{HS},
among others; and for hypergroups (as an algebraic generalization of
schemes) by French and Zieschang~\cite{FZ}. It has been analyzed
(evidently independently) in the setting of fusion categories and fusion
rings by Gelaki and Nikshych~\cite{GN}. They use
the term ``nilpotent'' for what the other authors above call ``residually
thin''. This usage has the desirable consequence that a finite group is
nilpotent in the classical sense if and only if its representation
category (resp. character ring) is nilpotent as a fusion category (resp.
fusion ring)~\cite [Remark~4.7]{GN}.
However, the group algebra of any finite group, as a fusion ring, is nilpotent according to their
definition. This seems to leave room for an alternative definition of
nilpotent. The one presented in this paper (see Definition~\ref{def:res_thin_nilp} below) is a direct generalization of the one
given for association schemes by Hanaki~\cite{H}. 

Our context here is the family of \emph{table algebras}, finite
dimensional algebras over the complex numbers with a distinguished basis
that satisfies certain axioms (see Definition~\ref{def:TA} below). The
adjacency (or Bose-Mesner) algebras of association schemes, group
algebras, and Hecke (double coset) algebras constructed from group
algebras are examples, and fusion rings comprise a subfamily.
Hypergroups, in the sense of~\cite{FZ}, are generalizations of table
algebras. Our main results include a formula for the degrees of the
basis elements of a residually thin table algebra, and a proof that this
formula holds only for such algebras (see Theorem~\ref{thm:thin-by-number} 
below). Gelaki and Nikshych's integrality theorem~\cite [Corollary~5.3]{GN} for fusion categories is an immediate consequence. These are evidently new results
for association schemes as well. We also define the notion of
\emph{thin central chains} for a table algebra, in particular upper,
lower, and bi-anchored such chains, and use it to exposit our alternate
definition of nilpotent table algebra. This definition coincides with the usual one
in the case of a finite group; and when applied to commutative table
algebras, it is equivalent to both the author's~\cite{Blau1} and Gelaki
and Nikshych's~\cite{GN} definitions. Theorems~\ref{thm:res_thin} and~\ref{thm:nilp} 
below relate residual depth to nilpotence class (see Definitions~\ref{def:depth}, \ref{def:class}), 
and thus extend 
Gelaki and
Nikshych's result for the commutative case~\cite [Theorem~4.16]{GN}.
Lower bi-anchored thin-central chains exist for any nilpotent table
algebra by Definition~\ref{def:ptc_chains}. Theorem~\ref{thm:nilp_lower} gives a partial characterization of them. But
unlike the case of the lower central series of a finite group, they are
not necessarily unique. Example~\ref{ex:lower} shows this, and thereby
gives a negative answer to a question of Hanaki~\cite [Question~2.11]{H} regarding uniqueness of lower central series in an association
scheme.
 

We recall a few well known definitions and facts needed in order to
state the main results. These are developed in a number of sources, in
particular~\cite{AFM, BZ,Blau2}.


 
\begin{defi}

\label{def:TA}
A \emph{table algebra} (TA) $(A, B)$ is a finite dimensional
algebra $A$
over the complex numbers $\C$, and a distinguished basis $B$ that
contains $1_A$, such that the following properties hold:
\begin{enumerate}[label=(1\alph*)]
\item\label{defi1.1_1a} %[(1a)]
The structure constants for $B$ are all nonnegative real numbers; that is, for all
$b,c \in B$,
\[
bc = \sum_{d \in B} \lambda_{bcd}d, \quad \text{for some }
\lambda_{bcd} \in \R_{\geq 0}.
\]
\item\label{defi1.1_1b} %[(1b)]
There is an algebra anti-automorphism (denoted by $^*$) of $A$ such
that $(a^*)^* = a$ for all $a \in A$; and $B^* = B$.
\item\label{defi1.1_1c} %[(1c)]
For all $b,c \in B$, 
$\lambda_{bc1} = 0$ if and only if $c \neq b^*$. 
\end{enumerate}
\end{defi}

%\

It follows as a consequence of the definition that $\lambda_{bb^*1} =
\lambda_{b^*b1} > 0$ for all $b \in B$. Frobenius-Perron eigenvalue
theory yields that for each table algebra there exists a unique algebra
homomorphism $\delta: A \rightarrow {\C}$, called the \emph{degree map},
such that $\delta(b) = \delta(b^*) > 0$ for all $b \in B$. The table
algebra $(A, B)$ is called \emph{standard} if for all $b \in B$,
$\delta(b) = \lambda_{bb^*1}$. Any table algebra can be \emph{rescaled}
(replace each $b \in B$ by $\beta_bb$, for suitable $\beta_b > 0$) to
one that is standard.

%\

Let $(A, B)$ be a standard table algebra (STA). For any subsets $S, T$
of $B$, the \emph{set product} $ST \coloneqq \cup_{s\in S, t\in
T}\Supp _B(st)$, $S^* \coloneqq \{s^* | s \in S\}$, and $S^+ \coloneqq \sum_{s
\in S}s$. Note that set product is associative. If $T = \{t\}$, a
singleton set, then $S\{t\}$ (resp. $\{t\}S, S\{t\}S$) is denoted $St$
(resp. $tS, StS$). The \emph{order} of a subset $S$ is $o(S) \coloneqq \delta
(S^+)$.

A nonempty subset $C \subseteq B$ is called \emph{closed} if $C^*C
\subseteq C$. In this case, $({\C}C, C)$ is again a table algebra, and
the set of \emph{left cosets} $Cb$ for $b \in B$ partition $B$, as do
the \emph{right cosets} $bC$, and the $C$-$C$ \emph{double cosets}
$CbC$. A \emph{quotient element}, for any $b \in B$, is $b//C \coloneqq 
o(C)^{-1}(CbC)^+$. Let $B//C \coloneqq \{b//C | b \in B\}$, and
$A//C \coloneqq {\C}(B//C)$. Then $(A//C, B//C)$ is a
STA, called the \emph{quotient algebra}, or \emph{double coset algebra}
of $(A, B)$ by the closed subset $C$. Its degree map is
$\delta\downarrow _{A//C}$, and its anti-automorphism is
${^*}\downarrow_{A//C}$. Furthermore, $o(B//C) =
o(B)/o(C)$. The closed subsets $D$ with $C \subseteq D \subseteq B$ are
in bijection with the closed subsets of $B//C$ via $D
\leftrightarrow D//C$ (see Proposition~\ref{prop:subset_corresp} below). The closed subset $C$ is called \emph{normal} (resp.
\emph{strongly normal}) in $B$ if $bC = Cb$ (resp. $bCb^* = C$) for all
$b \in B$. Strongly normal closed subsets are normal, but the converse
is not always true. 

An element $x \in B$ is called \emph{thin} (or \emph{linear}, or
\emph{grouplike}) if $xx^* = 1$. This is equivalent to $\delta(x) = 1$; 
and if $x$ is thin, then $xb \in B$ and $bx \in B$ for all $b \in B$.
So $\delta(xb) = \delta(bx) = \delta(b)$ for all $b \in B$. Now
$O_{\vartheta}(B)\coloneqq \{x \in B | xx^* = 1\} = \{x \in B |
\delta(x) = 1\}$ is a group, called the \emph{thin radical} of $B$; and
$B$ is called thin if $B = O_{\vartheta}(B)$. Obviously,
$O_{\vartheta}(B)$ is the unique maximal closed subset $D$ of $B$
such that $D$ is thin. For $C$ a closed subset of $B$, the quotient
$B//C$ is thin if and only if $C$ is strongly normal in $B$. Since the
intersection of strongly normal closed subsets is again strongly normal,
there is a unique minimal closed subset of $B$, denoted
$O^{\vartheta}(B)$, such that $B//O^{\vartheta}(B)$ is thin. Now
$O^{\vartheta}(B)$ is called the \emph{thin residue} of $B$, and it
equals the closed subset of $B$ generated by $\Supp _B(bb^*)$ for
all $b \in B$. (See~\cite[Theorem~2.3.1]{Z1} or~\cite [Theorem~3.2.1]{Z2}, 
where the algebraic proof for association schemes holds
verbatim for table algebras.) 

%\

Throughout, $Z(A)$ will denote the center of the algebra $A$, and $Z(B)$
will mean $B \cap Z(A)$.

\begin{defi}
\label{def:chain}
A \emph{chain $\mathcal{C}$ of length} $n$ is a collection $\{C_i\}_{i=0}^n$
of closed subsets of $B$ such that either $C_0 \subset C_1 \subset
\cdots \subset C_{n-1} \subset C_n$ or $C_n \subset C_{n-1} \subset
\cdots \subset C_1 \subset C_0$. 
 It is called a \emph{bi-anchored chain}
(\emph{B-chain}) if $C_0 = \{1\}$ and $C_n = B$ (or $C_n = \{1\}$ and
$C_0 = B$). It is a \emph{thin chain} (\emph{T-chain}) if $C_{i+1}//C_i$
is thin for all $0 \leq i \leq n-1$ (or $C_{i-1}//C_i$ is thin for all
$1 \leq i \leq n$). It is called a \emph{thin-central chain}
(\emph{TC-chain}) if $C_{i+1}//C_i$ is thin and $C_{i+1}//C_i \subseteq
Z(B//C_i)$ for all $0 \leq i \leq n-1$ (or $C_{i-1}//C_i$ is thin and
$C_{i-1}//C_i \subseteq Z(B//C_i)$ for all $1 \leq i \leq n$).
\end{defi}

\begin{defi}
\label{def:res_thin_nilp}
A STA $(A, B)$ is \emph{residually thin} if there exists a bi-anchored
thin chain (BT-chain). It is \emph{nilpotent} if there exists a
bi-anchored thin-central chain (BTC-chain).
\end{defi}

\begin{rema}
\label{rem:nilp_group}
It is immediate from the definition that for any finite group $G$, the
group algebra $({\C}G, G)$ is nilpotent as a STA if and only if $G$ is
nilpotent in the usual group-theoretic sense.
\end{rema}

\begin{prop}
\label{prop:res_thin_order}
Let $(A, B)$ be a residually thin STA, with BT-chain $B = C_0 \supset
C_1 \supset \cdots \supset C_{n-1} \supset C_n = \{1\}$. Then $o(C_i)$
is an integer, and $o(C_{i+1}) \vert o(C_i)$ for all $0 \leq i \leq
n-1$.
\end{prop}

\begin{proof}
Since each $C_j//C_{j+1}$ is a group, $o(C_j//C_{j+1}) =
o(C_j)/o(C_{j+1})$ is an integer for $0 \leq j \leq n-1$. Since $o(C_i)
= o(C_{i+1})o(C_i//C_{i+1}) = \prod_{j \geq i}o(C_j//C_{j+1})$, the
result follows.
\end{proof}

\begin{defi}
\label{def:card_number}
Let $(A, B)$ be a STA with a B-chain $\mathcal{C}$: $B = C_0 \supset C_1
\supset \cdots C_{n-1} \supset C_n = \{1\}$. For any $1 \leq i \leq n$
and any $b \in
B\backslash C_i$, define the positive integer $m(\mathcal{C}, i,
b)$ for $i < n$ by 
\[
m(\mathcal{C}, i, b) \coloneqq \prod_{j=i}^{n-1}\card (C_j//C_{j+1}
\cdot b//C_{j+1} \cdot C_j//C_{j+1}),
\]
where $C_j//C_{j+1} \cdot b//C_{j+1} \cdot C_j//C_{j+1}$ is the
$C_j//C_{j+1}$ - $C_j//C_{j+1}$ double coset of $b//C_{j+1}$ in
$B//C_{j+1}$; and $m(\mathcal{C}, n, b) \coloneqq \card (C_nbC_n) = \card\{b\} = 1$.
\end{defi}
 
We now can state our first main results.

\begin{thm}
\label{thm:thin-by-number}
Let $(A, B)$ be a STA with a B-chain $\mathcal{C}$: $B = C_0 \supset C_1
\supset \cdots \supset C_{n-1} \supset C_n = \{1\}$. Then $\mathcal{C}$
is a thin chain if and only if for all $i$ with $1 \leq i \leq n$ and
all $b \in C_{i-1}\backslash C_i$,
\[
\delta(b) = o(C_i)/m(\mathcal{C}, i, b).
\]
\end{thm}

\begin{thm}
\label{thm:thin-div}
Let $(A, B)$ be a STA with a BT-chain $\mathcal{C}$: $B = C_0 \supset
C_1 \supset \cdots \supset C_{n-1} \supset C_n = \{1\}$. Let $b \in
C_i$ for $i < n$. Then the following hold:
\begin{enumerate}[label = (\roman*)] 
\item $o(C_{i+1})/\delta(b)$ is an integer divisor of
$o(C_{i+1})^2$. In particular,
$o(O^{\vartheta}(B))/\delta(b)$ is an integer divisor of
$o(O^{\vartheta}(B))^2$ for all $b \in B$. 
\item If each $C_i$ is normal in $B$ (for example, if $\mathcal{C}$ is a
BTC-chain), then $o(C_{i+1})/\delta(b)$ is an integer divisor of
$o(C_{i+1})$; hence, $\delta(b)$ is its complementary integer divisor.
\end{enumerate}
\end{thm} 

Theorem~\ref{thm:thin-by-number} is proved in Section~\ref{sect3}  below, and
Theorem~\ref{thm:thin-div} in Section~\ref{sect4}. 

\begin{rema}
\label{rem:fus}
A fusion ring $(A, B)$ is a table algebra with integer structure
constants, and where $\lambda_{bb^*1} = 1$ for all $b \in B$. If
$\phi(b)$ is the degree of such $b$, then $\phi(b)^2$ is the degree of
the corresponding basis element in the rescaled standard basis. So
Gelaki and Nikshych's result~\cite [Corollary~5.3]{GN} that
$o(O^{\vartheta}(B))/\phi(b)^2$ is an integer follows immediately.
\end{rema}

\begin{defi}
\label{def:terminal}
Let $(A, B)$ be a STA. A TC-chain $\{Q_i\}_{i=0}^q$ is called
\emph{terminal} if $B = Q_0 \supset Q_1 \supset \cdots \supset Q_{q-1}
\supset Q_q$, and there is no closed subset $Q_{q+1} \subset Q_q$ such
that $\{Q_i\}_{i=0}^{q+1}$ is a TC-chain.
\end{defi}

\begin{defi}
\label{def:ptc_chains} 
Let $(A, B)$ be a STA.
\begin{enumerate}[label = (\roman*)] 
\item\label{defi1.11_i} The \emph{residual thin 
chain} $
\mathcal{R}$ is the chain $R_0 = B$, $R_1 =
O^{\vartheta}(B)$, $R_i = O^{\vartheta}(R_{i-1}) \mbox{ for }1 \leq i$. 
Thus, $\mathcal{R}$ has length $r$, where $r$ is the least nonnegative
integer with $R_{r+1} = R_r$.
\item\label{defi1.11_ii} The \emph{radical thin chain} $\mathcal{J}$ is the chain $J_0 = \{1\}, J_1 =
O_{\vartheta}(B), J_i//J_{i-1} = O_{\vartheta}(B//J_{i-1})$ for $1
\leq i$. Thus, $\mathcal{J}$ has length $j$, the least nonnegative
integer with $J_{j+1} = J_j$.
\item\label{defi1.11_iii} The \emph{upper thin-central chain} $\mathcal{Z}$ is the chain $Z_0 = \{1\}$, $
Z_1 = O_{\vartheta}(B) \cap Z(B), Z_i//Z_{i-1} =
O_{\vartheta}(B//Z_{i-1}) \cap Z(B//Z_{i-1}) \mbox{ for }1 \leq i$. Then
$\mathcal{Z}$ has length $u$, the least nonnegative integer with
$Z_{u+1} = Z_u$.
\item\label{defi1.11_iv} A \emph{lower thin-central chain of length} $q$ is a terminal TC-chain $\mathcal{Q}:
Q_0 \supset Q_1 \supset \cdots \supset Q_q$ 
such that $o(Q_q)$ is minimal, and $\sum_{i=0}^qo(Q_i)$ is 
minimal over all such chains. In particular, a lower \mbox{BTC-}chain
 of length $q$ is a BTC-chain $B = Q_0 \supset Q_1 \supset \cdots \supset
Q_q = \{1\}$ such that $\sum_{i=0}^qo(Q_i)$ is minimal. 
\end{enumerate}
\end{defi}

\begin{defi}
\label{def:abel_res}
Let $(A, B)$ be a STA. Define $O^{\alpha}(B)$, the \emph{thin abelian
residue}, to be the unique closed subset of $B$ with $B \supseteq
O^{\alpha}(B) \supseteq O^{\vartheta}(B)$ and
$O^{\alpha}(B)//O^{\vartheta}(B)$ is the commutator subgroup of the group
$B//O^{\vartheta}(B)$, that is, $[B//O^{\vartheta}(B),
B//O^{\vartheta}(B)]$. 
\end{defi}
 
\begin{rema}\ \relax
\label{rem:not_unique}

\begin{enumerate}[label = (\roman*)] 
\item\label{rema1.13_i} Proposition~\ref{prop:subset_corresp}\ref{prop2.5_ii} below shows that
$O^{\alpha}(B)$ is well-defined. From its definition and elementary
group theory, it is the unique smallest closed subset of
$B$ such that the quotient is an abelian group. 
\item\label{rema1.13_ii} If $(A, B)$ is the group algebra of a finite group $G$, then
$\{G_i\}_{i=0}^q$, where $G_0 = G, G_1 = [G, G], G_i = [G, G_{i-1}]$ for all
$1 \leq i \leq q$, where $G_{q+1} = G_q$, is the unique lower TC-chain
for $(A, B)$.
\item\label{rema1.13_iii} Lower TC-chains for an arbitrary STA $(A, B)$ always exist, by the
definition. 
If $(A, B)$ is nilpotent, then every lower BTC-chain of minimal
length 
begins $B = Q_0 \supset Q_1 = O^{\alpha}(B)$, but is not necessarily
unique. 
See Example~\ref{ex:lower} and Theorem 
\ref{thm:nilp_lower} below. The example yields a negative answer to 
a 
question of Hanaki~\cite [Question~2.11]{H} for association
schemes. 
\end{enumerate}
\end{rema} 

The following theorem is proved by French and Zieschang~\cite [Theorem~6.1]{FZ} in their more
general context of hypergroups. We include a short proof in Section~\ref{sect4},
for completeness. 

\begin{thm}
\label{thm:res_thin}
Let $(A, B)$ be a STA. If $\mathcal{C} = \{C_i\}_{i=0}^n$ is a T-chain,
$B = C_0 \supset C_1 \supset \cdots \supset C_n$, then $C_i \supseteq
R_i$ for all $0 \leq i \leq n$. Hence, $(A, B)$ is residually thin if
and only if the residual T-chain $\mathcal{R}$ is bi-anchored.
Furthermore, if $(A, B)$ is residually thin, and if $\mathcal{C}$ is any
BT-chain of length $n$, then $n \geq r$ \textup{(}$r =$ the length of
$\mathcal{R}$\textup{)}.
\end{thm}

\begin{defi} 
\label{def:depth}
The \emph{residual depth} of a residually thin STA $(A, B)$ is the
minimum length of all BT-chains in $B$; by Theorem~\ref{thm:res_thin},
it equals $r$, the length of $\mathcal{R}$.
\end{defi}

\begin{rema}
\label{rem:res_v_rad}
The radical T-chain $\mathcal{J}$ in a residually thin STA of depth $r$
may be bi-anchored of arbitrarily large length $j > r$, or it may not be
bi-anchored at all. See Example~\ref{ex:rad_length} below.
\end{rema}

\begin{defi}
\label{def_p-standard}
\cite [Definition~1.4]{BC} Let $p$ be a prime. A STA $(A, B)$ is called
\emph{$p$-standard} if $o(B) = p^N$ for some integer $N > 0$, and for
all $b \in B$, $\delta(b) = p^{n_b}$ for some integer $n_b \geq 0$. An
association scheme is called a \emph{$p$-scheme} if its adjacency
algebra is $p$-standard.
\end{defi}

\begin{prop}
\label{prop:p-stan_res_thin}
Let $(A, B)$ be a STA with $o(B) = p^N$ for some prime $p$ and integer
$N > 0$. Then $(A, B)$ is $p$-standard if and only if it is residually
thin.
\end{prop}

\begin{rema}
\label{rem:p-standard}
Not every $p$-standard table algebra is nilpotent. The example from
\cite [Remark~1.4]{BC}, which is taken from~\cite [No.10851]{HM}, is
$2$-standard of order 32, but has no nontrivial thin central basis elements.
\end{rema} 
 
We have the following analog of Theorem~\ref{thm:res_thin} for
thin-central chains and nilpotent STAs. It is a straightforward
generalization of~\cite [Theorem~2.5]{H}. 

\begin{thm}
\label{thm:nilp}
Let $(A, B)$ be a STA. If $\mathcal{C} = \{C_i\}_{i=0}^n$, $\{1\} = C_0
\subset C_1 \subset \cdots \subset C_n$ is a TC-chain, then $C_i
\subseteq Z_i$ for all $0 \leq i \leq n$. Hence, $(A, B)$ is nilpotent
if and only if the upper TC-chain $\mathcal{Z}$ is bi-anchored.
Furthermore, if $(A, B)$ is nilpotent, and if $\mathcal{C}$ is any
BTC-chain of length $n$, then $n \geq u$ \textup{(}$u =$ the length of
$\mathcal{Z}$\textup{)}.
\end{thm}

\begin{defi}
\label{def:class} 
The \emph{nilpotence class} of a nilpotent STA $(A, B)$ is the minimum
length of all the BTC-chains in $B$. By Theorem~\ref{thm:nilp}, it
equals $u$, the length of $\mathcal{Z}$.
\end{defi} 

 
It follows from Theorem~\ref{thm:res_thin} that for any nilpotent STA,
the nilpotence class is at least the residual depth. If $A$ is
commutative, then every T-chain is a TC-chain, hence by Theorem
\ref{thm:nilp} the residual depth
is at least the nilpotence class. So the following result of Gelaki and
Nikshych is immediate.

\begin{coro}
\label{cor:class_v_depth} 
\cite [Theorem~4.16]{GN} Let $(A, B)$ be a commutative STA. Then $(A,
B)$ is residually thin iff $\mathcal{R}$ is bi-anchored iff
$\mathcal{Z}$ is bi-anchored iff $(A, B)$ is nilpotent; and in this
case, the residual depth and nilpotence class of $(A, B)$ are equal.
\end{coro}

\begin{thm}
\label{thm:nilp_lower}
If $(A, B)$ is a nilpotent STA, then every lower BTC-chain of length $u$
begins $B \supset O^{\alpha}(B)$.
\end{thm}



Preliminary results (mostly known) and some further definitions are
collected in Section~\ref{sect2}. Theorem~\ref{thm:thin-by-number} and a related
more general result are proved in Section~\ref{sect3}. Section~\ref{sect4} contains proofs of Theorem~\ref{thm:res_thin}, Theorem~\ref{thm:thin-div},
Proposition~\ref{prop:p-stan_res_thin}, and other
structural results for residually thin STAs. 
 In particular, we
show that if $D$ is any closed subset of a residually thin table basis
$B$, then $o(D)$ is an integer such that $o(D) \mid o(B)$. Example~\ref{ex:rad_length} is also presented.
Section~\ref{sect5} establishes Theorem~\ref{thm:nilp} and Theorem~\ref{thm:nilp_lower}, and studies further aspects of TC-chains and nilpotent STAs. In particular, Corollary~\ref{cor:nilp_group} 
shows that a STA $(A, B)$ is nilpotent if and only if
$B//O^{\vartheta}(B)$ is nilpotent as a group, and $Z_k \supseteq
O^{\vartheta}(B)$ for some integer $k$. Example~\ref{ex:lower}
demonstrates 
that lower BTC-chains in a nilpotent STA are not always
unique. 



\section{Preliminaries}\label{sect2}
The results in the section for which proofs are omitted are known;
proofs for them are given 
in~\cite{AFM}, \cite{BZ}, or~\cite{Blau2}. Throughout,
$(A, B)$ is a table algebra (TA).



There is a positive definite Hermitian form $(\ ,\ )$ on $A$ such that
for
all $b,c,d \in B$,
\begin{center}
$(b,c) = \beta_{bc^*1}$; and
$(bc,d) = (b,dc^*) = (c,b^*d)$.
\end{center}

\begin{defi}
\label{def:TA-hom}
Let $(A, B), (U, V)$ be TAs. A \emph{table algebra} 
\textup{(}TA-\textup{)}\emph{homomorphism} 
$\psi: (A, B) \rightarrow (U, V)$ 
is an algebra homomorphism $\psi: A \rightarrow U, (\psi (1_A) = 1_U)$,
such that for each $b \in B$, there is some $v \in V$ and $\alpha_{b}
\in \R_{>0}$ with
$ \psi(b) = \alpha_{b}v$.
Define
$V_{\psi} \coloneqq \{ v \in V \mid \psi(b) = \alpha_bv \mbox{ for some }b \in
B \mbox{ and } \alpha_b \in \R_{>0}\}$. Then $V_{\psi}$ is a closed
subset of $V$. 
Now $\psi$ is called a TA-\emph{isomorphism} if $\psi$ is one-to-one
and $V_{\psi} = V$. This is equivalent to $\psi: A \rightarrow U$
being an algebra isomorphism. If there exists such a TA-isomorphism, 
we 
write $(A, B) \cong (U, V)$.
\end{defi}

\begin{defi}
\label{def:kernel}
The \emph{kernel} of a TA-homomorphism $\psi$ is defined as
$\ker \psi \coloneqq \{b \in B \vert \psi(b) = \alpha_b1_U \mbox{ for some
}\alpha_b \in \R_{>0}\}$.
\end{defi}

Then $\psi$ preserves the respective anti-automorphisms, $\ker \psi$ is a normal closed subset of $B$, and the
following ``Fundamental Homomorphism Theorem'' holds:
 

\begin{prop}
\label{prop:FHT}
Let $\psi : (A, B) \rightarrow (U, V)$ be a TA-epimorphism of STAs. Let
$C = \ker \psi$. Let $e = e_C \coloneqq o(C)^{-1}C^+$, a central
idempotent of $A$. Then $\pi : (A, B) \rightarrow (A//C, B//C)$, where
$\pi(b) = be$ for all $b \in B$, is a TA-epimorphism, as $be =
(\delta(b)/\delta(b//C))b//C$. Furthermore, there is a TA-isomorphism
$\overline{\psi} : (A//C, B//C) \rightarrow (U, V)$ such that
$\overline{\psi}\circ \pi = \psi$. Finally, since both $(A//C, B//C)$
and $(U, V)$ are standard, then for all $b \in B$,
$\overline{\psi}(b//C) = v$, where $\psi(b) = \alpha_bv$ for some
$\alpha_b \in {\R}_{>0}$. 
\end{prop} 


Proposition~\ref{prop:FHT} has the following consequence:

\begin{lemm}
\label{lem:isom}
Let $(A, B)$ be a STA with closed subsets $C, D$ where $D$ is normal in
$B$. Then $DC = CD$ is a closed subset, $D\cap C$ is a normal closed
subset of $C$, and
\[
C//(D\cap C) \cong CD//D,
\] 
via a TA-isomorphism that yields the correspondence $c//D\cap C
\leftrightarrow c//D$, for all $c \in C$.
\end{lemm}


The first two parts of the next proposition are contained in~\cite
[Proposition 2.13]{BZ}, and part~\ref{prop2.5_iii} is proved for hypergroups in~\cite
[Lemma 4.6]{FZ}. Parts~\ref{prop2.5_iii}, \ref{prop2.5_iv}, and~\ref{prop2.5_v} are proved below.

\begin{prop}
\label{prop:subset_corresp}
Let $(A, B)$ be a STA and $C$ a closed subset of $B$. Then the
following hold:
\begin{enumerate}[label = \textup{(}\roman*\textup{)}] 
\item\label{prop2.5_i} The correspondence $D \mapsto D//C$ is a bijection between the set of
closed subsets of $B$ that contain $C$ and the set of closed subsets of
$B//C$.
\item\label{prop2.5_ii} Suppose that $D$ is a closed subset of $B$ with $C \subseteq D \subseteq
B$. Then for all $b \in B$, $(b//C)//(D//C) = b//D$. Hence,
$(B//C)//(D//C) = B//D$.
\item\label{prop2.5_iii} Suppose that $C \subseteq D \subseteq B$. Then $D$ is strongly normal
in $B$ if and only if $D//C$ is strongly normal in $B//C$.
\item\label{prop2.5_iv} Suppose that $C \subseteq D \subseteq B$ and $D$ is normal in $B$. Then
$D//C$ is normal in $B//C$.
\item\label{prop2.5_v} Suppose that $C \subseteq D \subseteq B$ and $C$ is normal in $B$. Then
$D$ is a normal subset of $B$ if and only if $D//C$ is a normal subset
of $B//C$.
\end{enumerate}
\end{prop}

\begin{proof}
\ref{prop2.5_iii}~$D$ is strongly normal in $B$ iff $B//D$ is thin iff
$(B//C)//(D//C)$ is thin (by~\ref{prop2.5_ii}) iff $D//C$ is strongly normal in
$B//C$. 
 
\ref{prop2.5_iv}~Since $D$ is assumed normal in $B$, $(CbC)D = D(CbC)$ for all $b \in
B$. It follows that $b//C\cdot D//C = D//C\cdot b//C$, hence $D//C$ is
a normal subset of $B//C$.
 
\ref{prop2.5_v}~Suppose that $C$ is a normal subset of $B$, and that $D//C$ is
normal in $B//C$. Then for all $b \in B$, $b//C\cdot D//C = D//C\cdot
b//C$, hence $(CbC)D = D(CbC)$. But $C \subseteq D$, so $CD = D = DC$.
Thus we have $CbD = DbC$. Since $C$ is normal in $B$, $CbD = bCD = bD$,
and $DbC = DCb = Db$. The result follows.
\end{proof}


\begin{lemm}
\label{lem:sub_strongly}
Let $(A, B)$ be a STA with closed subsets $C_1 \supseteq C_2$ and $D$
such that 
$C_2$ is strongly normal in $C_1$, and for all $c \in C_1$, $cD
\subseteq DcC_2$. 
 Then $DC_1$ and
$DC_2$ are closed subsets such that $DC_2$ is
strongly normal in $DC_1$.
\end{lemm} 

\begin{proof} 
It follows from the hypothesis that $C_1D \subseteq DC_1$. Then $C_1D =
(DC_1)^* \supseteq (C_1D)^* = DC_1$. So $DC_1 = C_1D$, hence $DC_1$ is
a closed subset. 
The hypothesis restricted to $c \in C_2$ similarly yields $DC_2 = C_2D$,
hence $DC_2$ is closed. 
Any $b \in DC_1$ is in $\Supp(dc)$ for some $d \in D, c \in C_1$. Then
set products 
\[
bDC_2b^* \subseteq dcDC_2c^*d^* = d(cD)C_2c^*d^* \subseteq
d(DcC_2)C_2c^*d^* = D(cC_2c^*)d^*.
\] 
Now $C_2$ strongly normal in $C_1$ and $C_2D = DC_2$ yield
\[
D(cC_2c^*)d^* = DC_2d^* = D(C_2d^*) \subseteq D(DC_2) = DC_2.\qedhere
\]
\end{proof}

\begin{rema}
\label{rem:hypoth}
Given closed subsets $C_1 \supseteq C_2$ and $D$ with $C_2$ strongly
normal in $C_1$, the final hypothesis of Lemma~\ref{lem:sub_strongly}
will follow from either $D$ normal in $B$, or both $C_1$ and $C_2$
normal in $B$ with $(c//C_2)\cdot(DC_2//C_2) =
(DC_2//C_2)\cdot(c//C_2)$ for all $c \in C_1$. This is because the latter assumption
implies that $cDC_2 = (C_2cC_2)(DC_2) = (DC_2)(C_2cC_2) = DcC_2$.
\end{rema}
 
\begin{lemm}
\label{lem:intersect}
Let $(A, B)$ be a STA with closed subset $C$ and normal closed subset
$D$. If $O^{\vartheta}(C) \subseteq D$, then $C//(D\cap C)$ is thin.
\end{lemm}

\begin{proof}
Lemma~\ref{lem:sub_strongly} implies that $DO^{\vartheta}(C)$ is
strongly normal in $DC$. But $O^{\vartheta}(C) \subseteq D$ by
hypothesis, so $D$ is strongly normal in $DC$. Hence, $DC//D$ is thin.
By Lemma~\ref{lem:isom}, $C//(D\cap C)$ is thin.
\end{proof}

\begin{lemm}
\label{lem:constant}
Let $(A, B)$ be a STA, $C$ a closed subset of $B$, and $b \in B$. Then
$\delta$ is constant over the double coset $CbC$ if and only if $o(CbC)
= \card(CbC)\delta(b')$ for all $b' \in CbC$.
\end{lemm}

\begin{proof}
Since $o(CbC) = \sum_{b' \in CbC}\delta(b')$ and $\card(CbC)$ are
independent of any particular $b' \in CbC$, the proof is immediate.
\end{proof}

\begin{lemm}
\label{lem:action}
Let $(A, B)$ be a STA, $C$ a thin closed subset of $B$, and $b \in B$.
Then $\card (CbC) \mid o(C)^2$. If $C$ is normal in $B$, then
$\card(CbC) \mid o(C)$.
\end{lemm}

\begin{proof}
The group $C\times C$ acts on the double coset $CbC$ as follows: for
all $c_1, c_2, x, y, \in C$,
\[
(c_1bc_2)^{(x, y)} = x^{-1}c_1bc_2y.
\]
The double coset itself is the sole orbit under this action. Hence, if
$S \coloneqq \{(x, y) \in C\times C \mid x^{-1}by = b\}$, i.e. $S$ is the
stabilizer of $b$ under the action, then $\card (CbC) = |C\times C
: S|$. If $C$ is normal in $B$, then $CbC = bC$, and $C$ acts on $bC$ by
right multiplication. If $S^r$ denotes the stabilizer of $b$ in $C$
under this action, then $\card (CbC) = |C : S^r|$.
\end{proof}
 
\begin{lemm}
\label{lem:m_sq}
Let $(A, B)$ be a STA, and $\mathcal{C} = \{C_i\}_{i=0}^n$ a BT-chain
with $B = C_0 \supset C_1 \supset \cdots \supset C_n = \{1\}$. Then for
all $1 \leq i \leq n$ and all $b \in B\backslash C_i$, $m(\mathcal{C},
i, b) \mid o(C_i)^2$. If each $C_i$ is normal in $B$ \textup{(}in
particular, if $\mathcal{C}$ is a BTC-chain\textup{)}, then $m(\mathcal{C}, i, b)
\mid o(C_i)$.
\end{lemm}

\begin{proof}
This is immediate from the definitions if $i = n$, so assume $1 \leq i
< n$ and $b \in B\backslash C_i$. Lemma~\ref{lem:action} implies that
for each $j$ with $i \leq j \leq n-1$, 
\[
 \card (C_j//C_{j+1}\cdot b//C_{j+1}\cdot C_j//C_{j+1}) \mid
o(C_j//C_{j+1})^2 = o(C_j)^2/o(C_{j+1})^2.
\]
Hence,
\[
m(\mathcal{C}, i, b) = \prod_{j=i}^{n-1}\card (C_j//C_{j+1}\cdot
b//C_{j+1}\cdot C_j//C_{j+1}) \mathrel{\bigg|}
\prod_{j=i}^{n-1}\frac{o(C_j)^2}{o(C_{j+1})^2} = o(C_i)^2.
\]
If each $C_i$ is normal in $B$, then each $C_j//C_{j+1}$ is a thin
normal closed subset of $B//C_{j+1}$, by Proposition
\ref{prop:subset_corresp}\ref{prop2.5_iv}. So Lemma~\ref{lem:action} yields that 
for all $i \leq j \leq n-1$, 
\[
\card (C_j//C_{j+1}\cdot b//C_{j+1}\cdot C_j//C_{j+1}) =
\card (b//C_{j+1}\cdot C_j//C_{j+1}) \mid o(C_j//C_{j+1}).
\]
Therefore, $m(\mathcal{C}, i, b) \mid
\prod_{j=i}^{n-1}(o(C_j)/o(C_{j+1})) = o(C_i)$.
\end{proof}


 
\section{Degrees and double cosets}\label{sect3} 

Throughout this section, $(A, B)$ is a STA with degree map $\delta$, and
$\mathcal{C} = \{C_i\}_{i=0}^n$ is a B-chain with $B = C_0 \supset C_1
\cdots \supset C_n = \{1\}$. We prove a theorem that, for a given
B-chain, yields a criterion whereby the degree of a quotient element is
found in terms of the cardinality of double cosets in quotient bases
determined by the chain. The result is applied to prove Theorem
\ref{thm:thin-by-number}.

\begin{defi}
\label{def:quochain}
Fix integer $k$ with $0 < k < n$. Define the chain
\[
\mathcal{C}//C_k \coloneqq \{C_i//C_k\}_{i=0}^k.
\]
\end{defi}

\begin{rema}
\label{rem:quochain}
We have by Proposition~\ref{prop:subset_corresp}\ref{prop2.5_i} that
\[
B//C_k = C_0//C_k \supset C_1//C_k \cdots \supset C_{k-1}//C_k \supset
C_k//C_k,
\]
so that $\mathcal{C}//C_k$ is indeed a B-chain of length $k$ in
$B//C_k$.
\end{rema}

\begin{lemm}
\label{lem:chain_num_mult}
Fix integers $i, k$ with $0 \leq i \leq k \leq n$, so that $C_i
\supseteq C_k$.
Then for all $b \in B\backslash C_i$,
\[
m(\mathcal{C}, i, b) = m(\mathcal{C}//C_k, i, b//C_k)\cdot
m(\mathcal{C}, k, b).
\]
\end{lemm}
 
\begin{proof}
If $i = k$, then $m(\mathcal{C}//C_k, k, b//C_k) = 1$ by definition, as
$\mathcal{C}//C_k$ has length $k$. If $k = n$, then $m(\mathcal{C}, n,
b) = 1$, $\mathcal{C}//C_n = \mathcal{C}$, and $b//C_n = b$. So we may
assume that $i < k < n$. 
Proposition~\ref{prop:subset_corresp}\ref{prop2.5_ii} implies for all $i \leq j
\leq k-1$, 
\[
\left((C_j//C_k)//(C_{j+1}//C_k)\right) \cdot
\left((b//C_k)//(C_{j+1}//C_k)\right) \cdot
\left((C_j//C_k)//(C_{j+1}//C_k)\right)
\] 
\[
 = C_j//C_{j+1}\cdot b//C_{j+1}\cdot
C_j//C_{j+1},
\]
so that by Definition~\ref{def:card_number}, $m(\mathcal{C}//C_k, i,
b//C_k) = \prod_{i \leq j \leq k-1}\card (C_j//C_{j+1}\cdot
b//C_{j+1}\cdot C_j//C_{j+1})$. Since by the same definition $m(\mathcal{C}, k, b)$ (resp.
$m(\mathcal{C}, i, b)$) equals the analogous product over $k \leq j \leq
n-1$ (resp. $i \leq j \leq n-1$), the result follows.
\end{proof}

\begin{thm}
\label{thm:constant_2coset}
Let $\mathcal{C} = \{C_i\}_{i=0}^n$ be a B-chain with $B = C_0 \supset
C_1 \supset \cdots \supset C_n = \{1\}$. Then $\delta$ is constant on
each double coset $C_j//C_{j+1}\cdot b//C_{j+1}\cdot C_j//C_{j+1}$ for
all $0 \leq i \leq j \leq n-1$ and all $b \in B\backslash C_i$ if and only if,
for all $0 \leq i \leq n-1$,
\[
\delta(b//C_i) = \frac{\delta(b)}{o(C_i)}m(\mathcal{C}, i, b).
\]
\end{thm}
 
\begin{proof}
Fix $i \geq 0$. By Proposition~\ref{prop:subset_corresp}\ref{prop2.5_ii}, $\delta$ is constant on
$C_j//C_{j+1}\cdot b//C_{j+1}\cdot C_j//C_{j+1}$ for all $i \leq j
\leq n-1$ and all $b \in B\backslash C_i$ if and only if $\delta$ is
constant on 
\begin{multline}\label{eq:mod_n-1} 
\left((C_j//C_{n-1})//(C_{j+1}//C_{n-1})\right)\cdot
\left((b//C_{n-1})//(C_{j+1}//C_{n-1})\right)
\\
\cdot
\left((C_j//C_{n-1})//(C_{j+1}//C_{n-1})\right)
%\begin{equation}
%\\
\mbox{ for all }i \leq j \leq n-2 \mbox{ and all }b \in
B\backslash C_i, 
\end{multline}
and
\begin{equation}
\label{eq:just_n-1}
\delta \mbox{ is constant on } C_{n-1}bC_{n-1} \mbox{ for all }b \in
B\backslash C_i.
\end{equation} 
 
Since $B\backslash C_i$ is a union of $C_{n-1}$-$C_{n-1}$ double cosets,
Lemma~\ref{lem:constant} implies that~\eqref{eq:just_n-1} is equivalent
to $\delta(b//C_{n-1}) = \frac{\delta(b)}{o(C_{n-1})}m(\mathcal{C}, n-1,
b)$ for all $b \in B\backslash C_i$. So if $i = n-1$, the theorem is
proved. 

Suppose that $i < n-1$. By induction on $n$, \eqref{eq:mod_n-1}~is
equivalent to 
\[
\delta\left((b//C_{n-1})//(C_i//C_{n-1})\right) = 
\frac{\delta(b//C_{n-1})}{o(C_i//C_{n-1})}m(\mathcal{C}//C_{n-1}, i,
b//C_{n-1})
\]
for all $b \in B\backslash C_i$. So if~\eqref{eq:mod_n-1} and~\eqref{eq:just_n-1} both hold, we have
\[
\begin{aligned}
\delta(b//C_i) &= \delta \left((b//C_{n-1})//(C_i//C_{n-1})\right) =
\frac{\delta(b//C_{n-1})}{o(C_i)/o(C_{n-1})}m(\mathcal{C}//C_{n-1}, i,
b//C_{n-1})
\\
&= \frac{\delta(b)}{o(C_{n-1})o(C_i)/o(C_{n-1})}m(\mathcal{C}, n-1,
b)\cdot m(\mathcal{C}//C_{n-1}, i, b//C_{n-1}).
\end{aligned}
\]
Hence by Lemma~\ref{lem:chain_num_mult}, $\delta(b//C_i) =
\frac{\delta(b)}{o(C_i)}m(\mathcal{C}, i, b)$.

Conversely, suppose that for all $0 \leq i \leq n$ and all $b \in
B\backslash C_i$, $\delta(b//C_i) =
\frac{\delta(b)}{o(C_i)}m(\mathcal{C}, i, b)$. If $i = n-1$, we have
already shown that $\delta$ is constant on $C_{n-1}bC_{n-1}$. Suppose
that $i < n-1$. Then by Lemma~\ref{lem:chain_num_mult}, 
\[
\begin{aligned}
\delta(b//C_i) &= \frac{\delta(b)}{o(C_i)}m(\mathcal{C}, n-1, b) \cdot
m(\mathcal{C}//C_{n-1}, i, b//C_{n-1})
\\
 &= \frac{o(C_{n-1})}{o(C_i)}\delta(b//C_{n-1})m(\mathcal{C}//C_{n-1},
i, b//C_{n-1}).
\end{aligned}
\]
So 
\[
\delta\left((b//C_{n-1})//(C_i//C_{n-1})\right) =
\frac{\delta(b//C_{n-1})}{o(C_i//C_{n-1})}m(\mathcal{C}//C_{n-1}, i,
b//C_{n-1})
\] 
for all $0 \leq i \leq n-2$ and $b \in B\backslash C_i$, whence
$\mathcal{C}//C_{n-1}$ satisfies the same hypothesis as~$\mathcal{C}$.
Then by induction on $n$, $\delta$ is constant on all double cosets
\begin{multline*}
\left((C_j//C_{n-1})//(C_{j+1}//C_{n-1})\right)\cdot
\left((b//C_{n-1})//(C_{j+1}//C_{n-1})\right)
\\
 \cdot
\left((C_j//C_{n-1})//(C_{j+1}//C_{n-1})\right) = C_j//C_{j+1} \cdot
b//C_{j+1}\cdot C_j//C_{j+1},
\end{multline*} 
for all $i \leq j \leq n-2$ and $b \in
B\backslash C_i$. This establishes the converse. 
 \end{proof}

\begin{proof} [Proof of Theorem~\ref{thm:thin-by-number}]

Observe that $C_{n-1}$ is thin iff $\delta(b) = 1$ for all $b \in
C_{n-1}\backslash C_n$ iff $\delta(b) =
o(C_n)/m(\mathcal{C}, n, b)$.

If $\mathcal{C}$ is thin, then $C_j//C_{j+1}$ is a thin closed subset
for all $0 \leq j \leq n-1$, by definition. So $\delta$ is constant on
each double coset $C_j//C_{j+1}\cdot b//C_{j+1}\cdot C_j//C_{j+1}$ for
all $0 \leq i \leq j \leq n-1$ and all $b \in B\backslash C_i$. Hence, $\delta(b//C_i) =
\frac{\delta(b)}{o(C_i)}m(\mathcal{C}, i, b)$ for all $0 \leq i \leq
n-1$ and $b \in B\backslash C_i$, by Theorem~\ref{thm:constant_2coset}.
But $C_{i-1}//C_i$ thin and $b \in C_{i-1}\backslash C_i$ imply that
$\delta(b//C_i) = 1$. Hence, $\delta(b) = o(C_i)/m(\mathcal{C},
i, b)$ for $0 \leq i \leq n-1$.

Conversely, suppose that $\delta(b) = o(C_i)/m(\mathcal{C}, i, b)$ for
all $0 \leq i \leq n$ and all $b \in C_{i-1}\backslash C_i$. In
particular if $i = n$, then $C_{n-1}$ is thin, as noted above. So
$\delta$ is constant on $C_{n-1}bC_{n-1}$ for all $b \in B$.

If $i < n$ and $b \in C_{i-1}\backslash C_i$, then $C_{n-1}$ thin, Lemma~\ref{lem:constant} and Lemma
\ref{lem:chain_num_mult} imply that
\[
\delta(b//C_{n-1}) \Mk = \Mk  \frac{\delta(b)}{o(C_{n-1})}m(\mathcal{C}, n-1,
b) \Mk =\Mk  \frac{o(C_i)m(\mathcal{C}, n-1, b)}{m(\mathcal{C}, i, b)o(C_{n-1})}
\Mk =\Mk  \frac{o(C_i//C_{n-1})}{m(\mathcal{C}//C_{n-1}, i, b//C_{n-1})}.
\]
So the same hypothesis holds for the chain $\mathcal{C}//C_{n-1}$ in
$B//C_{n-1}$ as for $\mathcal{C}$. Induction on $n$ implies that
$\mathcal{C}//C_{n-1}$ is thin. Since $C_j//C_{j+1} =
(C_j//C_{n-1})//(C_{j+1}//C_{n-1})$ for all $0 \leq j \leq n-2$ by
Proposition~\ref{prop:subset_corresp}\ref{prop2.5_ii}, $\mathcal{C}$ is thin.
\end{proof}


\section{Residually thin STAs}\label{sect4}


In this section, we prove Theorem~\ref{thm:res_thin}, Theorem~\ref{thm:thin-div}, Proposition~\ref{prop:p-stan_res_thin}, and in
Theorem~\ref{thm:res_thin_sub_quo} 
properties of residually thin STAs. 
 Parts of the latter result are proved for
hypergroups in~\cite {FZ}, as noted below; but we include proofs here,
as they fit easily into our context. Example~\ref{ex:rad_length}
below illustrates how much residual depth and radical length can differ.
Throughout, $(A, B)$ denotes a STA.

\begin{proof} [Proof of Theorem~\ref{thm:res_thin}] 
Let $\mathcal{C} = \{C_i\}_{i=0}^n$ be a T-chain in $B$, $B = C_0
\supset C_1 \supset \cdots \supset C_{n-1} \supset C_n$. Then $B//C_1$
is thin, by Definition~\ref{def:chain}. Since $R_1 = O^{\vartheta}(B)$
is the unique minimal closed subset of $B$ such that $B//R_1$ is thin,
$C_1 \supseteq R_1$. Suppose that $C_i \supseteq R_i$ for some $i \geq
1$. Since by definition $C_i//C_{i+1}$ is thin, $C_{i+1}$ is (strongly)
normal in $C_i$. Hence, $R_iC_{i+1} = C_{i+1}R_i$ is a closed subset of
$C_i$. So by Proposition
\ref{prop:subset_corresp}\ref{prop2.5_i}, $R_iC_{i+1}//C_{i+1}$ is a subgroup of
$C_i//C_{i+1}$. By Lemma~\ref{lem:isom}, $R_i//(R_i\cap C_{i+1}) \cong
R_iC_{i+1}//C_{i+1}$. Thus, $R_i//(R_i\cap C_{i+1})$ is thin. Then by
the definition of $\mathcal{R}$, $R_{i+1} \subseteq R_i\cap C_{i+1}
\subseteq C_{i+1}$. It follows by induction that $C_i \supseteq R_i$
for $0 \leq i \leq n$. So if $C_n = \{1\}$, then $R_n = \{1\}$ and $n
\geq r$. The theorem follows.
\end{proof}

\begin{proof} [Proof of Theorem~\ref{thm:thin-div}]
Let $b \in C_i$. We may assume that $b \neq 1$. So for some $j$ with
$i \leq j < n$, $b \in C_j\backslash C_{j+1}$. Theorem
\ref{thm:thin-by-number} then implies that $o(C_{j+1})/\delta(b) =
m(\mathcal{C}, j+1, b)$, which is an integer. Lemma~\ref{lem:m_sq}
yields further that $o(C_{j+1})/\delta(b) \mid o(C_{j+1})^2$. As $j+1 \leq i+1$,
$o(C_{j+1}) \mid o(C_{i+1})$ by Proposition~\ref{prop:res_thin_order}.
Hence, $o(C_{i+1})/\delta(b)$ is an integer and $o(C_{i+1})/\delta(b) 
\mid o(C_{i+1})^2$. If each $C_j$ is normal in $B$, then Lemma
\ref{lem:m_sq} implies that $o(C_{j+1})/\delta(b) \mid o(C_{j+1})$, so
$o(C_{i+1})/\delta(b) \mid o(C_{i+1})$. 

By Theorem
\ref{thm:res_thin}, the residual thin chain $\mathcal{R}$ is also
bi-anchored. Since $b \in B = R_0$, and $R_1 = O^{\vartheta}(B)$, our
proof shows that $o(O^{\vartheta}(B))/\delta(b)$ is an integer divisor
of $o(O^{\vartheta}(B))^2$.
\end{proof}

\begin{proof} [Proof of Proposition~\ref{prop:p-stan_res_thin}]
Suppose that $(A, B)$ is $p$-standard. By~\cite [Proposition~3.2]{BC},
there is a chain $\{1\} \subset C_1 \subset C_2 \subset \cdots \subset
C_N = B$ with $C_{i+1}//C_i$ a cyclic group of order $p$ for $0 \leq i <
N$. So by Definition~\ref{def:res_thin_nilp}, $(A, B)$ is residually
thin.
Suppose that $(A, B)$ is residually thin. Let $\{R_i\}_{i=0}^r$ be the
residual thin chain, so that $B = R_0 \supset R_1 \supset \cdots \supset
R_r = \{1\}$, by Theorem~\ref{thm:res_thin}. By Proposition
\ref{prop:res_thin_order}, $o(R_1)$ is an integer divisor of $o(B) =
p^N$. Thus, $o(R_1) = p^{N_1}$ for some nonnegative integer $N_1 < N$.
For each $b \in B$, $p^{N_1}/\delta(b) = o(R_1)/\delta(b)$ is an integer
divisor of $o(R_1)^2 = p^{2N_1}$ by Theorem~\ref{thm:thin-div}. It
follows that $\delta(b)$ is a power of $p$, and $(A, B)$ is
$p$-standard.
\end{proof}
 
Part~\ref{theo4.1_i}  and most of~\ref{theo4.1_ii} of the next result are proved for hypergroups by
French and Zieschang in
\cite [Theorem~6.3]{FZ}. Since the proofs in our context are 
shorter, they are included below. Part~\ref{theo4.1_iii} seems new.

\begin{thm}
\label{thm:res_thin_sub_quo}
Let $(A, B)$ be a residually thin STA of depth $n$, with a BT-chain
$\{C_i\}_{i=0}^n$, $B = C_0 \supset C_1 \supset \cdots \supset C_n =
\{1\}$. Let $D$ be any closed subset of $B$.
\begin{enumerate}[label = \textup{(}\roman*\textup{)}] 
\item\label{theo4.1_i} The distinct members of $\{D\cap C_i\}_{i=0}^n$ form a BT-chain for $D$,
hence $D$ is residually thin of depth at most $n$. In particular,
$o(D)$ is an integer.
\item\label{theo4.1_ii} Suppose either that $D$ is normal in $B$, or that each $C_i$ is normal in $B$
and $(c//C_{i+1})\cdot(DC_{i+1}//C_{i+1}) =
(DC_{i+1}//C_{i+1})\cdot(c//C_{i+1})$ for all $c \in C_i$. Then the distinct members of
$\{DC_i//D\}_{i=0}^n$ form a BT-chain for $B//D$, hence $B//D$ is
residually thin of depth at most $n$.
\item\label{theo4.1_iii} $o(D) \mid o(B)$.
\end{enumerate}
\end{thm}

\begin{proof}
Let $b \in D\cap C_i$. Then $b(D\cap C_{i+1})b^* \subseteq C_{i+1}$,
since $C_{i+1}$ is strongly normal in $C_i$; and $b(D\cap C_{i+1})b^* \subseteq
D$, since $b \in D$ and $D$ is closed. Thus, $b(D\cap C_{i+1})b^*
\subseteq D\cap C_{i+1}$. So $D\cap C_{i+1}$ is strongly normal in $D\cap
C_i$. Hence, $(D\cap C_i)//(D\cap C_{i+1})$ is thin. So $D$ is
residually thin of depth at most $n$, and $o(D)$ is an integer by
Proposition~\ref{prop:res_thin_order}. Thus, \ref{theo4.1_i}~follows.

Assume the hypotheses of~\ref{theo4.1_ii}. Then each $DC_i$ is closed, and the
hypotheses of Lemma~\ref{lem:sub_strongly} hold for each pair $C_i$,
$C_{i+1}$, by Remark~\ref{rem:hypoth}. By 
Lemma~\ref{lem:sub_strongly}, each $DC_{i+1}$ is strongly normal in
$DC_i$. So $DC_i//DC_{i+1} = (DC_i//D)//(DC_{i+1}//D)$ is thin.
Therefore, the distinct members of $\{DC_i//D\}_{i=0}^n$ form a BT-chain
for $B//D$, and~\ref{theo4.1_ii} is proved.

Now $C_1$ is residually thin of depth $n-1$, and is strongly normal in
$B$; and $D\cap C_1$ is a closed
subset of $C_1$. So $o(D\cap C_1) \mid o(C_1)$ by induction on
$n$. Since $DC_1//C_1$ is a subgroup of the group $B//C_1$, we have
$o(DC_1//C_1) \mid o(B//C_1) = o(B)/o(C_1)$. By Lemma
\ref{lem:isom}, $D//(D\cap C_1) \cong DC_1//C_1$. Thus,
\[
o(D)/o(D\cap C_1) = o(D//(D\cap C_1)) = o(DC_1//C_1) \mid 
o(B)/o(C_1).
\]
Hence, 
\[
o(D) \mid \frac{o(B)o(D\cap C_1)}{o(C_1)} =
\frac{o(B)}{o(C_1)/o(D\cap C_1)} \mid o(B),
\]
where the last division holds because $o(C_1) \mid o(B)$ by
Proposition~\ref{prop:res_thin_order}, and $o(D\cap C_1) \mid 
o(C_1)$. Therefore, \ref{theo4.1_iii}~holds. 
\end{proof}

\begin{exam}
\label{ex:rad_length}
Let $G$ be a finite group, $H$ a proper subgroup of $G$, $Y$ any other
finite group, and $C = H\times Y$. Define the group homomorphism $\psi
: C \rightarrow H$ by $\psi(hy) = h$ for all $h \in H$, $y \in Y$.
Consider the partial wreath product $C \circ_{\psi}G = C\cup
(G\backslash H)$ as a basis for a vector space over $\C$~\cite
[Definition 4.1, Definition 4.2]{Blau1}, \cite [Definition~1.13,
Definition~2.4]{Blau3}. That is, $H$ is replaced inside $G$ by $C$.
Multiplication in $\C(C\circ_{\psi}G)$ is defined as follows: for all
$h \in H, y \in Y, g \in G\backslash H$, $(hy)\cdot g = hg$ (the product
in $G$), and $g\cdot (hy) = gh$. If $g_1, g_2 \in G\backslash H$ with
$g_1g_2 = h \in H$, then in $\C(C\circ_{\psi}G)$, $g_1g_2 =
\frac{1}{o(Y)}hY^+ = \frac{1}{o(Y)}Y^+h$. The other products are as in
$G$ or $C$. Rescale each $g \in G\backslash H$ as $b \coloneqq o(Y)g$. Define
$B$ as the rescaled $C\circ_{\psi}G$ and $A \coloneqq {\C}B$. Then $(A, B)$ is a
STA. For all $c \in C$, $c^* = c^{-1}$ and $cc^* = 1$. For all $g \in
G\backslash H$ and $b = o(Y)g$, $b^* = o(Y)g^{-1}$ and $bb^* = o(Y)Y^+$.
Hence, $O_{\vartheta}(B) = C$, and $O^{\vartheta}(B) = Y$. Then the
residual thin chain $\mathcal{R}$ is
\[
 B \supset O^{\vartheta}(B) = Y \supset \{1\},
\]
and $(A, B)$ is residually thin of depth 2.

Suppose that $H = N_G(H)$; that is, $H$ is its own normalizer in the
group $G$. Then for all $b = o(Y)g \in B\backslash C$, the set product
\[
bCb^* = \{o(Y)g_1 \mid g_1 \in gHg^{-1}\backslash H\} \cup
((gHg^{-1}\cap H)\times Y) 
\not\subseteq H \times Y,
\] 
since $g \notin N_G(H)$. So if $D$ is any closed subset of $B$ with $C
\subset D \subseteq B$, then $C$ is not strongly normal in $D$, hence
$D//C$ is not thin. It follows that the radical thin chain
$\mathcal{J}$ is $\{1\} \subset C$, of length 1 and not bi-anchored.

Let $n > 1$ be an arbitrary integer, and suppose that $G$ is a group with a chain of subgroups $H = H_1 < H_2 <
\cdots < H_n = G$, such that $N_G(H_i) = H_{i+1}$ for $1 \leq i \leq
n-1$. (This is the case, for example, if $G = D_{2^n}$, the dihedral
group of order $2^n$, $u$ is a noncentral involution in $G$, $\langle v
\rangle$ is the cyclic subgroup of order $2^{n-1}$, and $H_i = \langle
v^{2^{n-i}}\rangle \cup \langle v^{2^{n-i}}\rangle u$ for $1 \leq i \leq
n$.) Let $C_0 = \{1\}, C_1 = C$, and $C_i = C\cup \{o(Y)g \mid g \in
H_i\backslash H\}$ for $2 \leq i \leq n$. Then for each $1 \leq i \leq n-1$, $C_i$
is a closed subset that is strongly normal in $C_{i+1}$, but $bC_ib^*
\not\subseteq C_i$ for any $b \in B\backslash C_{i+1}$. It follows that
$C_{i+1}//C_i = O_{\vartheta}(B//C_i)$. Hence, $\{C_i\}_{i=0}^n$ is the
radical thin chain $\mathcal{J}$ of $B$, and it is bi-anchored with
length $n$. 
 
\end{exam}

\begin{rema}
\label{rem:rad_length} 
Regard the groups $G$, $C$ above as association
schemes in the usual way. (For $g \in G$, the relation $g_L$ on
underlying set $G$ is given by $(x, y) \in g_L$ iff $xy^{-1} = g$ for
all $x, y \in G$; and similarly for $C$.) The adjacency algebras are
isomorphic as table algebras to the group algebras ${\C}G$ and ${\C}C$,
and the partial wreath product constructed above is realized as the
adjacency algebra of the wedge product of the schemes~\cite [Section~3]{M}. 
 So Example~\ref{ex:rad_length} applies to
association schemes.
\end{rema}



\section{Nilpotent STAs}\label{sect5}

We turn now to thin-central chains and nilpotent STAs. Theorems
\ref{thm:nilp}, \ref{thm:closed_sub_nilp}, one direction of Theorem~\ref{thm:quotient_nilp}, and
their proofs given below, are generalized directly from Hanaki's results
in~\cite {H}. Again, 
$(A, B)$ is always a STA.

\begin{proof} [Proof of Theorem~\ref{thm:nilp}]
Let $\{C_i\}_{i=0}^n$, $\{1\} \subset C_1 \subset \cdots \subset C_{n-1}
\subset C_n$ be a TC-chain. Note first that $C_1$ is normal in $B$; and
then by Proposition~\ref{prop:subset_corresp}\ref{prop2.5_v} all $C_i$ are normal
in $B$. Now $C_1 = C_1//C_0$ is thin, and $C_1 \subseteq Z(B)$. So
$C_1 \subseteq O_{\vartheta}(B)\cap Z(B) = Z_1$. 

Suppose that $C_i \subseteq Z_i$ for some $i \geq 1$. Now $C_i$ is
strongly normal in $C_{i+1}$, and $Z_i$ is normal in $B$. Hence, Remark
\ref{rem:hypoth} and Lemma~\ref{lem:sub_strongly} imply that $Z_i =
Z_iC_i$ is strongly normal in $Z_iC_{i+1}$. Thus, $Z_iC_{i+1}//Z_i$ is thin.
Since $C_{i+1}//C_i \subseteq Z(B//C_i)$, for any $c \in C_{i+1}$ and $b
\in B$, $c//C_i\cdot b//C_i = b//C_i\cdot c//C_i$. Hence, $cbC_i = bcC_i$, which,
as $C_i \subseteq Z_i$, implies that $cbZ_i = bcZ_i$. Since $c//Z_i$ is
thin, $c//Z_i\cdot b//Z_i = y//Z_i$ for any $y \in cbZ_i$, and similarly
for $b//Z_i\cdot c//Z_i$. Therefore, 
$c//Z_i\cdot b//Z_i = b//Z_i\cdot c//Z_i$, whence $c//Z_i \in Z(B//Z_i)$. Then
$Z_iC_{i+1}//Z_i \subseteq O_{\vartheta}(B//Z_i)\cap Z(B//Z_i) =
Z_{i+1}//Z_i$, and so $C_{i+1} \subseteq Z_{i+1}$. The theorem follows.
\end{proof} 

\begin{lemm}
\label{lem:to_lower}
Let $(A, B)$ be a STA, and let $\mathcal{C} =
\{C_i\}_{i=0}^q$, $C_0 \supset C_1 \supset \cdots \supset C_q$ 
be a TC-chain of length $q$ such that $C_q$ is a normal closed subset of
$B$. 
 Let $Q$ be a closed subset of $B$ such that $B \supseteq Q \supseteq
O^{\vartheta}(B)$ and $Q//O^{\vartheta}(B)$ is normal in
$B//O^{\vartheta}(B)$. Define $Q_i \coloneqq Q\cap C_i$ for $0 \leq i \leq q$. 
Then the distinct members of $\{Q_i\}_{i=0}^q$ form a TC-chain in $B$ of length
at most $q$.
\end{lemm}

\begin{proof}
Each quotient $Q_i//Q_{i+1} = (Q\cap C_i)//(Q\cap C_{i+1})$
is thin, by Theorem~\ref{thm:res_thin_sub_quo}\ref{theo4.1_i}. Since
$B//O^{\vartheta}(B)$ is a group, our hypothesis that
$Q//O^{\vartheta}(B)$ is normal in $B//O^{\vartheta}(B)$ implies that it
is strongly normal. Hence, $Q$ is strongly normal in $B$ by Proposition~\ref{prop:subset_corresp}\ref{prop2.5_iii}. Each $C_i//C_{i+1}$ is normal in
$B//C_{i+1}$, since $C_i//C_{i+1} \subseteq Z(B//C_{i+1})$ by the
definition of a TC-chain. Since $C_q$ is normal in $B$, it follows from
Proposition~\ref{prop:subset_corresp}\ref{prop2.5_v} that each $C_i$ is normal in $B$.
 
 If $x \in Q_i$ and
$b \in B$, then $x//Q_{i+1}$ thin implies that $x//Q_{i+1}\cdot
b//Q_{i+1}\cdot x^*//Q_{i+1}$ is a basis element in $B//Q_{i+1}$. Since
$Q$ is strongly normal in $B$ and $x^* \in Q$, $\Supp (bx^*b^*)
\subseteq Q$, hence $\Supp (xbx^*b^*) \subseteq Q$. Also, $x
\in C_i$, and $C_i//C_{i+1} \subseteq Z(B//C_{i+1})$
imply that $x//C_{i+1}\cdot b//C_{i+1}\cdot x^*//C_{i+1} =
b//C_{i+1}$. Therefore, $\Supp (x//C_{i+1}\cdot
b//C_{i+1}\cdot x^*//C_{i+1}\cdot b^*//C_{i+1})$ contains
$1//C_{i+1}$. Since $C_{i+1}$ is normal in $B$, this says that
$\Supp (xbx^*b^*C_{i+1})\cap C_{i+1} \neq \emptyset$, hence (via the
Hermitian form), $\Supp (xbx^*b^*)\cap C_{i+1} \neq \emptyset$. Since
$Q_{i+1} = Q\cap C_{i+1}$, $\Supp (xbx^*b^*)\cap Q_{i+1} \neq
\emptyset$. So $1//Q_{i+1} \in \Supp (x//Q_{i+1}\cdot
b//Q_{i+1}\cdot x^*//Q_{i+1}\cdot b^*//Q_{i+1})$, and thus
$x//Q_{i+1}\cdot b//Q_{i+1} = b//Q_{i+1}\cdot x//Q_{i+1}$. Therefore,
$Q_i//Q_{i+1} \subseteq Z(B//Q_{i+1})$ for all $i$. 
Thus, the distinct terms of $\{Q_i\}_{i=0}^q$ form a
TC-chain of length at most $q$. 
\end{proof} 


\begin{proof} [Proof of Theorem~\ref{thm:nilp_lower}] 
Let $\mathcal{C} = \{C_i\}_{i=0}^u$ be a BTC-chain of length $u$, the
nilpotence class of $(A, B)$, such that $\sum_{i=0}^uo(C_i)$ is minimal.
Then $\mathcal{C}$ is immediately a lower BTC-chain. Write
$\mathcal{C}$ as $B = C_0 \supset C_1 \supset \cdots \supset C_u =
\{1\}$. Let $Q = O^{\alpha}(B)$. By Lemma~\ref{lem:to_lower}, the
distinct terms of $\{Q\cap C_i\}_{i=0}^u$ form a TC-chain of length at
most $u$. Now $B//C_1 \subseteq Z(B//C_1)$ and is thin, so $B//C_1$ is an
abelian group. Therefore, $C_1 \supseteq Q$. Hence, $Q\cap C_1 = Q =
Q\cap C_0$. So this TC-chain starts with $Q$ and has length at most
$u-1$. Since $B//Q$ is an abelian group, we have that the distinct
terms of $\{B\}\cup \{Q\cap C_i\}_{i=0}^u$ form a BTC-chain of
length at most $u$. But Theorem~\ref{thm:nilp} implies that the length
is at least $u$; so we have that all terms are distinct and form a
terminal BTC-chain of length $u$, 
\begin{equation}
\label{eq:lower}
B \supset Q = Q\cap C_1 \supset Q\cap C_2 \supset \cdots \supset Q\cap
C_{u-1} \supset Q\cap C_u = \{1\}. 
\end{equation}
Now $o(B) + \sum_{i=1}^uo(Q\cap C_i) \leq \sum_{i=0}^uo(C_i)$. But
minimality of the latter implies that the sums are equal. Since $Q\cap
C_i \subseteq C_i$, and hence $o(Q\cap C_i) \leq o(C_i)$ for $1 \leq i
\leq n$, it follows
that $Q\cap C_i = C_i$. Thus, $C_1 = Q\cap C_1 = Q$, and every lower
BTC-chain of length $u$ begins 
$B \supset O^{\alpha}(B)$. 
\end{proof}
 
 
\begin{thm}
\label{thm:closed_sub_nilp}
Let $(A, B)$ be a nilpotent STA of class $u$, and let $B = C_u \supset
C_{u-1} \supset \cdots \supset C_1 \supset C_0 = \{1\}$ be a BTC-chain
of length $u$. Let $V$ be any closed subset of $B$. Then the distinct
members of $\{V\cap C_i\}_{i=0}^u$ form a BTC-chain for $V$; hence
$({\C}V, V)$
is nilpotent of class at most $u$. 
\end{thm}

\begin{proof}
Each $C_i$ is normal in $B$, hence each $V\cap C_i$ is normal in $V$.
By Theorem~\ref{thm:res_thin_sub_quo}\ref{theo4.1_i}, $(V\cap C_{i+1})//(V\cap C_i)$
is thin for $0 \leq i \leq u-1$. Since $(V\cap C_{i+1})C_i//C_i
\subseteq C_{i+1}//C_i$, which is central in $B//C_i$, then $(V\cap
C_{i+1})C_i//C_i$ is central in $VC_i//C_i$. 

By Lemma~\ref{lem:isom}, $V//(V\cap C_i) \cong VC_i//C_i$, where the
isomorphic correspondence between table bases is $v//(V\cap C_i)
\leftrightarrow v//C_i$ for all $v \in V$. Under this bijection,
$(V\cap C_{i+1})//(V\cap C_i) \leftrightarrow (V\cap C_{i+1})C_i//C_i$.
It follows that $(V\cap C_{i+1})//(V\cap C_i)$ is central in $V//(V\cap
C_i)$. Therefore, the distinct members of $\{V\cap C_i\}_{i=0}^u$
comprise a BTC-chain for $V$; and $({\C}V, V)$ is nilpotent of class at
most $u$.
\end{proof} 

Recall that the upper TC-chain $\mathcal{Z} = \{Z_i\}$ is
defined for any STA $(A, B)$ in Definition~\ref{def:ptc_chains}; and is
bi-anchored if and only if $(A, B)$ is nilpotent, by Theorem
\ref{thm:nilp}.
 
\begin{thm}
\label{thm:quotient_nilp}
Let $(A, B)$ be a STA, and $D$ a closed subset of $B$. Then $(A,
B)$ is nilpotent if and only if $(A//D, B//D)$ is nilpotent and $Z_k
\supseteq D$ for some $k \geq 0$.
\end{thm}

\begin{proof}
Suppose that $(A, B)$ is nilpotent, say of class $u$, so that $Z_u = B
\supseteq D$. Since all the $Z_i$ are normal in $B$, Theorem 
\ref{thm:res_thin_sub_quo}\ref{theo4.1_i} implies that the distinct terms of
$\{DZ_i//D\}_{i=0}^u$
form a BT-chain
\[
B//D = DZ_u//D \supseteq DZ_{u-1}//D \supseteq \cdots \supseteq
DZ_1//D \supseteq DZ_0//D = D//D
\]
for $B//D$. Then for all $y \in Z_i$ and $b \in B$,
\begin{multline*}
y//Z_{i-1}\cdot b//Z_{i-1} = b//Z_{i-1}\cdot y//Z_{i-1} \Rightarrow
 ybZ_{i-1} = byZ_{i-1}
\\
\Rightarrow ybDZ_{i-1} = byDZ_{i-1}
 \Rightarrow DZ_{i-1}ybDZ_{i-1} = DZ_{i-1}byDZ_{i-1}.
\end{multline*} 
Now $Z_{i-1}$ normal in $B$ and $y//Z_{i-1}$ central in $B//Z_{i-1}$ imply
that $DZ_{i-1}y = Z_{i-1}yD = yDZ_{i-1}$. So
\[
DZ_{i-1}ybDZ_{i-1} = y(DZ_{i-1}bDZ_{i-1}) = 
(DZ_{i-1}yDZ_{i-1})(DZ_{i-1}bDZ_{i-1}).
\]
Similarly, $DZ_{i-1}byDZ_{i-1} = (DZ_{i-1}bDZ_{i-1})(DZ_{i-1}yDZ_{i-1})$.
Therefore,
\begin{equation} 
\label{eq:equal_sets} 
(DZ_{i-1}yDZ_{i-1})(DZ_{i-1}bDZ_{i-1}) = 
(DZ_{i-1}bDZ_{i-1})(DZ_{i-1}yDZ_{i-1}).
\end{equation}
Since $y//DZ_{i-1}$ is thin, both $y//DZ_{i-1}\cdot b//DZ_{i-1}$ and
$b//DZ_{i-1}\cdot y//DZ_{i-1}$ are single basis elements in
$B//DZ_{i-1}$. So it follows from~\eqref{eq:equal_sets} that
\[
y//DZ_{i-1}\cdot b//DZ_{i-1} = b//DZ_{i-1}\cdot yDZ_{i-1}.
\] 
Then by Proposition~\ref{prop:subset_corresp}\ref{prop2.5_ii}, 
\begin{multline*}
(y//D)//(DZ_{i-1}//D)\cdot (b//D)//(DZ_{i-1}//D) 
\\
= (b//D)//(DZ_{i-1}//D)\cdot (y//D)//(DZ_{i-1}//D).
\end{multline*}
Since every element of $DZ_{i}//D$ has the form $y//D$ for some $y \in
Z_i$, this shows that $(DZ_i//D)//(DZ_{i-1}//D) \subseteq Z((B//D)//(DZ_{i-1}//D))$.
Therefore, the distinct terms of $\{DZ_i//D\}_{i=0}^u$ form a BTC-chain
for $B//D$, so that $(A//D, B//D)$ is nilpotent.

Conversely, suppose that $(A//D, B//D)$ is nilpotent and $Z_k \supseteq
D$ for some $k \geq 0$. Then 
$(A//Z_k, B//Z_k) = ((A//D)//(Z_k//D), (B//D)//(Z_k//D))$ is nilpotent,
by Proposition~\ref{prop:subset_corresp}\ref{prop2.5_ii} and the first part of this
proof. Hence, there is a TC-chain
\[
Z_k//Z_k = Y_0//Z_k \subset Y_1//Z_k \subset \cdots \subset Y_m//Z_k =
B//Z_k.
\]
So $Y_i//Y_{i-1} = (Y_i//Z_k)//(Y_{i-1}//Z_k)$ is thin
and contained in 
\[
Z((B//Z_k)//(Y_{i-1}//Z_k)) = Z(B//Y_{i-1}),
\]
 for each
$1 \leq i \leq m$. Thus,
\[
\{1\} = Z_0 \subset Z_1 \subset \cdots \subset Z_k \subset Y_1 \subset
\cdots \subset Y_{m-1} \subset Y_m = B
\]
is a BTC-chain for $B$, and $(A, B)$ is nilpotent. 
\end{proof}
 
\begin{coro}
\label{cor:nilp_group}
A STA $(A, B)$ is nilpotent if and only if $B//O^{\vartheta}(B)$ is a
nilpotent group and $Z_k \supseteq O^{\vartheta}(B)$ for some integer $k
\geq 0$.
\end{coro}

\begin{exam}
\label{ex:lower}
Fix an odd prime $p$. Let $G$ be a $p$-group of nilpotence class $u
\geq 3$, hence with lower central series
\[
G = G_0 \supset G_1 = [G,G] \supset G_2 = [G, G_1] \supset \cdots
\supset G_{u-2} \supset G_{u-1} \supset G_u = \{1\}.
\]
Assume furthermore that $G_{u-2}$ is abelian, and $|G_{u-1}| = p$.
(This holds, for example, if $G$ is the multiplicative group of upper
unitriangular $(u+1)\times(u+1)$ ($u \geq 3$) matrices over the field
${\F}_p$.) Let $H = G_{u-2}$; $Y = \langle y \rangle$, a group of order
$p$; $C = H\times Y$, and $\psi: C \rightarrow H$ be the group
homomorphism $hy^s \mapsto h$ for all $h \in H$, $y^s \in Y$. As in
Example~\ref{ex:rad_length}, let $B$ be the standard rescaling of the
partial wreath product $C \circ_{\psi} G$, and $A = {\C}B$. Then
\[
B = \{b_g \coloneqq pg \mid g \in G\backslash H\} \cup C,
\]
where $b_g(hy^s) = b_gh$, $(hy^s)b_g = hb_g$ for all $g \in G\backslash
H$, $h \in H$, $y^s \in Y$; and if $g_1, g_2 \in G\backslash H$, then
\[
b_{g_1}b_{g_2} 
= \begin{cases}
phY^+ & \mbox{ if }g_1g_2 = h \in H,\\
pb_{g_1g_2} & \mbox{ if }g_1g_2 \notin H. 
\end{cases}
\] 
In particular, $y^sb_g = b_gy^s = b_g$, and $b_gb_{g^{-1}} = pY^+$, for
all $g \in G\backslash H$. Thus, $b_g^* = b_{g^{-1}}$. Now $B//Y \cong
G$ via the correspondence $b_g//Y \leftrightarrow g$ for all $g \in
G\backslash H$ and $hy^s//Y \leftrightarrow h$ for all $h \in H$, $y^s
\in Y$. We so identify the two groups.

Define 
\[
Q_i \coloneqq \{b_g \mid g \in G_i\backslash H\} \cup C, \ \ 0 \leq i \leq
u-2.
\]
Then $Q_i$ is a closed subset, $Q_i \supseteq Y$, and $Q_i//Y = G_i$.
So for $1 \leq i \leq u-2$,
\[
Q_{i-1}//Q_i = (Q_{i-1}//Y)//(Q_i//Y) =
G_{i-1}//G_i
\]
\begin{equation}
\label{eq:lcentral}
 \mbox{ is thin and is central in } \\ G//G_i = (B//Y)//(Q_i//Y)
= B//Q_i. 
\end{equation}
Because $\{G_i\}$ is the lower central series for $G$, each $Q_i$, for
$1 \leq i \leq u-2$, is the unique minimal closed subset in $Q_{i-1}$ 
among all closed subsets that contain $Y$ such that~\eqref{eq:lcentral}
holds. But any closed subset of $B$ that is not contained in $C$
contains $b_g$ for some $g \in G\backslash H$, and so contains
$\Supp(b_gb_g^*) = Y$. Therefore, for $1 \leq i \leq u-2$, $Q_i$
is in fact the unique minimal closed subset in $Q_{i-1}$ so that~\eqref{eq:lcentral} is true. Hence, $\sum_{i=0}^{u-2}o(Q_i)$ is
uniquely minimal
for all TC-chains of length $u-2$ that proceed down from $B$. 

Let $G_{u-1} = \langle z \rangle$, of order $p$ by our choice of $G$.
Fix any integer $j$ with $1 \leq j < p$, and let $D_j \coloneqq \langle z^jy
\rangle$. Then $G_{u-1}$ central in $G$ and $Y$ central in $B$ imply
that each $D_j \subseteq Z(B)$. Since $C$ is an abelian group (as $H = G_{u-2}$ is abelian,
again by choice of $G$), $Q_{u-2}//D_j = C//D_j$ is thin and abelian.
Furthermore, for all $x = hy^r \in Q_{u-2}$ and $b_g \in B\backslash
Q_{u-2}$, we have in the nilpotent group $G$ that $hg = ghz^t$ for some
integer $t$. So in $B$, $b_g = pg$, $y^rb_g = b_gy^r = b_g$, and $C$
abelian yield
\[
xb_g = hy^rb_g = hb_g = b_ghz^t = b_ghy^rz^t = b_gxz^t = b_gz^tx.
\]
Now $z^t = z^{js}$ for some integer $s$, and $b_gz^{js} = b_g(z^jy)^s$.
Hence,
\[
xb_gD_j = b_gx(z^jy)^sD_j = b_gxD_j.
\]
Therefore, $x//D_j\cdot b_g//D_j = b_g//D_j\cdot x//D_j$, thus
$Q_{u-2}//D_j \subseteq Z(B//D_j)$. Since $D_j \subseteq Z(B)$ and
$o(D_j) = p$, it follows that for each $1 \leq j < p$, 
\[
B \supset Q_1 \supset Q_2 \supset \cdots \supset Q_{u-2} \supset D_j
\supset \{1\}
\]
is a lower BTC-chain for $B$. Each $Q_i$ is the unique minimal closed
subset of $Q_{i-1}$ that exists in a BTC-chain, for all $i \leq u-2$;
but each $D_j$ is minimal such in $Q_{u-2}$. Thus there are $p-1$
stringently minimal such chains, not a unique one. 
As in Example~\ref{ex:rad_length}, the
algebraic construction here is realized as the adjacency algebra of the wedge product of
association schemes. Hence this example too applies to association
schemes. 
 So the answer to~\cite [Question~2.11]{H} is negative.
\end{exam}


Our final result displays the role played in general by the subsets
$\Supp(b^*b)$ for $b \in B$ in finding TC-chains in $B$ from the
top down, with each term minimal in the previous one.

\begin{thm}
\label{thm:minimal}
Let $U$ be a closed subset in a STA $(A, B)$. Let $V =
O^{\alpha}(U)$. Let $S$ be the closed subset of $U$ such that $S
\supseteq V$ and $S//V = \langle \Supp(b^*//V\cdot b//V)\cap
(U//V)
 \mid b \in B\rangle $. Assume that $U//V$ is normal in $B//V$. Then the following hold:
\begin{enumerate}[label = \textup{(}\roman*\textup{)}] 
\item There is a unique closed subset $C$ of $U$ that
is minimal \textup{(}with respect to inclusion\textup{)} such that $C \supseteq S$, $C//V$
is normal in $B//V$,
 and $U//C \subseteq Z(B//C)$.
\item If $D$ is any closed subset of $U$ that is minimal with respect to
inclusion such that $V \subseteq D$, $D//V$
is normal in $B//V$, and $U//D \subseteq Z(B//D)$, then 
$D \subseteq C$.
\end{enumerate}
\end{thm}

\begin{proof}
If $Y$ is any closed subset with $V \subseteq Y \subseteq U$, then
$U//V$ an abelian group implies that $Y//V$ is a normal abelian subgroup
of $U//V$. Hence by $V$ normal in $U$ and Proposition
\ref{prop:subset_corresp}\ref{prop2.5_v}, $Y$
is normal in $U$; and $U//Y$ is an abelian group, in particular is thin.

Suppose that $C$ is a closed subset of $U$ such that $V \subseteq C$,
$C//V$ is normal in $B//V$, and $U//C \subseteq Z(B//C)$. Then $U//C$ thin
implies that for any $x \in U$ and $b \in B$, 
$x//C\cdot b//C\cdot x^*//C = b//C$. Hence, 
\[
 (x//V)//(C//V)\cdot
(b//V)//(C//V)\cdot (x^*//V)(C//V) = (b//V)//(C//V),
\] 
which, since $C//V$ is thin and normal in $B//V$, implies that
\[
x//V\cdot b//V\cdot x^*//V = b//V\cdot c_{x,b}//V, \ \mbox{ some
}c_{x,b}
\in C.
\]
Suppose that $D$ is another closed subset of $U$ with $V \subseteq D$,
$D//V$
normal in $B//V$, and $U//D \subseteq Z(B//D)$. 
Since $V \subseteq C\cap D
\subseteq U$, $U//C\cap D$ is also thin. An argument similar to
the one above 
yields that for all $x \in U$ and $b \in B$,
\[
x//V\cdot b//V\cdot x^*//V = b//V\cdot d_{x,b}//V, \ \mbox{ some
}d_{x,b}
\in D.
\]
Therefore, $b//V\cdot c_{x,b}//V = b//V\cdot d_{x,b}//V$, so that $b//V\cdot
(c_{x,b}//V)\cdot (d_{x,b}//V)^{-1} = b//V$. Then the thin element
$(c_{x,b}//V)\cdot (d_{x,b}//V)^{-1}$ is in $\Supp(b^*//V\cdot b//V)$, by the Hermitian
form for the STA $B//V$. Thus, $(c_{x,b}//V)\cdot (d_{x,b}//V)^{-1} \in S//V$. 

Suppose that $C \supseteq S$. Then $d_{x,b}//V \in (C//V)\cdot
(c_{x,b}//V) 
= C//V$. Therefore, $d_{x,b} \in C\cap D$. Now $V \subseteq C\cap D$,
and both closed subsets are normal in $U$. Hence,
\begin{multline*}
x//V\cdot b//V\cdot x^*//V = b//V\cdot d_{x,b}//V\\ 
\begin{aligned}
&\Rightarrow 
Vxbx^*V = Vbd_{x,b}V
\\
 &\Rightarrow (C\cap D)xbx^*(C\cap D) = (C\cap
D)bd_{x,b}(C\cap D)
\\
& \Rightarrow\begin{aligned}[t](x//C\cap D)\cdot (b//C\cap D)\cdot
(x^*//C\cap D) &= (b//C\cap D)\cdot (d_{x,b}//C\cap D)\\
 &= b//C\cap D,
 \end{aligned}
\end{aligned}
\end{multline*}
since $d_{x,b} \in C\cap D$ implies that $d_{x,b}//C\cap D = 1//C\cap
D$. It follows that $U//(C\cap D) \subseteq Z(B//(C\cap D))$. 
Now 
 $C//V$ and $D//V$ are normal in $B//V$. 
 So if $x \in C\cap
D$ and $b \in B$, then 
\[
x//V\cdot b//V = b//V\cdot c//V = b//V\cdot d//V,
\]
for some $c \in C$, $d \in D$ (both of which depend on $x$ and $b$).
Hence, $b//V\cdot c//V\cdot (d//V)^{-1} = b//V$, so that 
\[
 c//V\cdot (d//V)^{-1} \in \Supp (b^*//V\cdot b//V)\cap (U//V) \subseteq S//V
\subseteq C//V.
\]
Therefore, $d//V \in C//V$, so $d \in C\cap D$. We now have that
$((C\cap D)//V)\cdot (b//V) \subseteq (b//V)\cdot ((C\cap D)//V)$. Replacing
$b$ by $b^*$ and then applying the anti-automorphism yields the opposite
containment, hence $(C\cap D)//V$ is normal in $B//V$. 

We have shown
that $C\cap D$ satisfies the same hypotheses as $D$, provided that $C
\supseteq S$. Both claims of the theorem follow immediately. 
\end{proof}

\begin{thebibliography}{99}
\bibitem{AFM}Arad, Zvi; Fisman, Elsa; Muzychuk, Mikhail Generalized table algebras, Israel J. Math., Volume 114 (1999), pp. 29-60
\bibitem{Blau1}Blau, Harvey I. Quotient structures in C-algebras, J. Algebra, Volume 177 (1995) no. 1, pp. 297-337
\bibitem{Blau2}Blau, Harvey I. Table algebras, European J. Combin., Volume 30 (2009) no. 6, pp. 1426-1455
\bibitem{Blau3}Blau, Harvey I. Fusion rings with few degrees, J. Algebra, Volume 396 (2013), pp. 220-271
\bibitem{BC}Blau, Harvey I.; Chen, Shengan Normal series and character values in p-standard table algebras, Comm. Algebra, Volume 45 (2017) no. 11, pp. 4646-4655
\bibitem{BZ}Blau, Harvey I.; Zieschang, Paul-Hermann Sylow theory for table algebras, fusion rule algebras, and hypergroups, J. Algebra, Volume 273 (2004) no. 2, pp. 551-570
\bibitem{FZ}French, Christopher; Zieschang, Paul-Hermann On residually thin hypergroups, J. Algebra, Volume 551 (2020), pp. 93-118
\bibitem{GN}Gelaki, Shlomo; Nikshych, Dmitri Nilpotent fusion categories, Adv. Math., Volume 217 (2008) no. 3, pp. 1053-1071
\bibitem{H}Hanaki, Akihide Nilpotent schemes and group-like schemes, J. Combin. Theory Ser. A, Volume 115 (2008) no. 2, pp. 226-236
\bibitem{HM}Hanaki, Akihide; Miyamoto, Izumi Classification of association schemes with small vertices (2021) (http://math.shinshu-u.ac.jp/\textasciitilde{}hanaki/as/)
\bibitem{HS}Hanaki, Akihide; Shimabukuro, Osamu Indecomposable decompositions of modular standard modules for two families of association schemes, J. Algebraic Combin., Volume 46 (2017) no. 2, pp. 445-453
\bibitem{M}Muzychuk, Mikhail A wedge product of association schemes, European J. Combin., Volume 30 (2009) no. 3, pp. 705-715
\bibitem{Z1}Zieschang, Paul-Hermann An algebraic approach to association schemes, Lecture Notes in Mathematics, 1628, Springer-Verlag, Berlin, 1996, xii+189 pages
\bibitem{Z2}Zieschang, Paul-Hermann Theory of association schemes, Springer Monographs in Mathematics, Springer-Verlag, Berlin, 2005, xvi+283 pages
\end{thebibliography}
\end{document}
